\section{Fuentes de datos}

Podemos dividir las distintas fuentes de datos en tres tipos, dependiendo de qué proceso vamos a seguir para su extracción.


\subsection*{\textit{Web scraping}}
\begin{itemize}
    \item ESA Sky Data Hub \cite{esa-sky-scraping}: Plataforma de datos espaciales que incluye imágenes del espacio profundo
    \item NASA Astronomy Picture of the Day (APOD) Random \cite{nasa-apod-scraping}: Acceso aleatorio a la colección de imágenes astronómicas diarias proporcionadas por la NASA.
\end{itemize}


\subsection*{Datos estáticos}
\begin{itemize}
  \item NUFORC (National UFO Reporting Center) Dataset
  \item SETI (Search for Extraterrestrial Intelligence) Data
  \item ESA (European Space Agency) Open Data
  \item SatNogs Dataset
\end{itemize}


\subsection*{APIs externas}
\begin{itemize}
  \item UFO Stalker API
  \item MUFON (Mutual UFO Network) API
  \item NASA open-access APIs
  \item NY2O API (rate-limited)
  \item Google Maps API
  \item Mapbox API
  \item NASA Image and Video Library API: 
  \item NASA Astronomy Picture of the Day (APOD) API \cite{nasa-apod-api}: Acceso selectivo a la colección de imágenes astronómicas diarias proporcionadas por la NASA.
  \item Sentinel Hub API \cite{sentinel-api}: Acceso selectivo a la colección de imágenes satélitales terrestres de los satélites Sentinel de la ESA
\end{itemize}